\documentclass[11pt]{article}
\usepackage{amsmath,amssymb,array,booktabs,longtable,geometry,fontspec,xcolor}
\usepackage{tikz}\usetikzlibrary{arrows.meta}
\usepackage[hidelinks]{hyperref}
\definecolor{claimlink}{RGB}{25,70,150}
\newcommand{\pvar}[2]{\hyperlink{#1-P#2}{\textcolor{claimlink}{P_{#2}}}}
\newcommand{\pvartarget}[2]{\raisebox{\ht\strutbox}{\hypertarget{#1-P#2}{}}}
\makeatletter
\newenvironment{derivation}
  {\begin{list}{}{\setlength{\leftmargin}{2.6em}\setlength{\labelwidth}{2.2em}%
    \setlength{\labelsep}{0.4em}\setlength{\itemsep}{2pt}\setlength{\parsep}{0pt}%
    \setlength{\topsep}{6pt}\@beginparpenalty\@M\@itempenalty\@M%
    \interlinepenalty\@M\raggedright}}
  {\end{list}}
\newcommand{\breakrow}[1]{\@itempenalty\z@\item[#1]\@itempenalty\@M}
\newcommand{\keepwithnext}{\@endparpenalty\@M}
\makeatother
\clubpenalty=10000
\widowpenalty=10000
\newcommand{\gloss}[1]{\enspace---\enspace{\small\itshape #1}}
\geometry{margin=1in}
\setlength{\emergencystretch}{3em}
\frenchspacing
\IfFontExistsTF{Times New Roman}
  {\setmainfont{Times New Roman}[Ligatures=TeX]}
  {\IfFontExistsTF{FreeSerif}{\setmainfont{FreeSerif}[Ligatures=TeX]}{}}
\IfFontExistsTF{DejaVu Sans Mono}
  {\setmonofont{DejaVu Sans Mono}[Scale=MatchLowercase]}
  {\IfFontExistsTF{FreeMono}{\setmonofont{FreeMono}[Scale=MatchLowercase]}{}}
\title{Born in Bethlehem — Micah 5:2}
\author{\credit{TestimonyArguments.BornInBethlehem}}
\date{Generated by \texttt{lake exe argtex}}
\newcommand{\credit}[1]{\ifcsname credit@#1\endcsname\csname credit@#1\endcsname\fi}
\newcommand{\creditline}[1]{\ifcsname credit@#1\endcsname{\centering\credit{#1}\par\medskip}\fi}
\InputIfFileExists{credits.tex}{}{}
\begin{document}
\maketitle
\noindent This is one argument of the Testimony library, printed on its own. The library prints every argument together in \texttt{arguments.pdf}. This one is set as its Lean source sets it: the module docstrings as prose, in the order they are written, and under each one the package, line of reason, countermodel or result it introduces. Atoms are numbered propositional variables, numbered once per argument, so a variable means the same claim in every position within it; the legend gives each one's claim, its classification, the tradition that holds it and the source that carries it.

\noindent Nothing here asserts that the premises are true --- only that the conclusion does or does not follow from them.

\tableofcontents
\clearpage
\section*{Born in Bethlehem — Micah 5:2}
\addcontentsline{toc}{section}{Born in Bethlehem — Micah 5:2}

\subsection*{TestimonyArguments.BornInBethlehem — Phase 1 vertical slice}
Matthew 2:5–6 quotes Micah 5:2 as grounds that the Messiah is born in Bethlehem. Both the Christian predictive reading and a critical reading are encoded, and the argument is carried through to a \texttt{Satisfaction\allowbreak{}Witness} — the first end-to-end connection from cited premises to a claim that a person meets a messianic criterion.

\texttt{jesus\allowbreak{}Born\allowbreak{}In\allowbreak{}Bethlehem} is cited to scripture alone. That is visible in the generated manifest via \texttt{scripture\allowbreak{}Only\allowbreak{}Atoms}, and it should be: the historicity of the Bethlehem birth is disputed in critical scholarship, and an argument that assumes it on the authority of the texts whose reading is in question is assuming part of what it sets out to show.

The argument is short enough to stay in one file. It runs on a \textbf{single line of reason}, which is exactly why it is the weakest of the fulfilment arguments and why \texttt{critical\_\allowbreak{}not\_\allowbreak{}establishes} costs the rival so little: deny the one interpretive premise and nothing is left standing. Compare \texttt{Born\allowbreak{}Of\allowbreak{}AVirgin}, which runs on four.

\subsubsection*{The claims, numbered}
\begin{longtable}{@{}r >{\raggedright\arraybackslash}p{0.30\textwidth} l >{\raggedright\arraybackslash}p{0.19\textwidth} >{\raggedright\arraybackslash}p{0.21\textwidth}@{}}
\toprule
 & claim & kind & held & source \\
\midrule
\endhead
\pvartarget{TestimonyArguments.BornInBethlehem}{1}$P_{1}$ & Micah 5:2 is a forward-looking Messianic prediction & interpretive & Christian, typological, well supported & \cite[ad loc. Mic 5:2]{keil-delitzsch-minor-prophets-1949} \\
\pvartarget{TestimonyArguments.BornInBethlehem}{2}$P_{2}$ & Matthew 2:6 quotes Micah 5:2 & textual & critical scholarship, consensus & \cite[app. Matt 2:6]{na28-2012}; \cite{ubs5-2014} \\
\pvartarget{TestimonyArguments.BornInBethlehem}{3}$P_{3}$ & Matthew’s quotation intends predictive fulfilment & interpretive & Christian, hist.-gramm., well supported & \cite[ad loc. Matt 2:6]{france-matthew-2007} \\
\pvartarget{TestimonyArguments.BornInBethlehem}{4}$P_{4}$ & Jesus of Nazareth was born in Bethlehem\textsuperscript{\dag} & historical & Christian, hist.-gramm., disputed & Matt 2:1; Luke 2:4-7 \\
\pvartarget{TestimonyArguments.BornInBethlehem}{5}$P_{5}$ & Micah 5:2 is a near-term oracle about a contemporary Judaean ruler & interpretive & critical scholarship, plausible & \cite{brown-birth-messiah-1993} \\
\pvartarget{TestimonyArguments.BornInBethlehem}{6}$P_{6}$ & The Messiah must be born in Bethlehem & interpretive & Christian, typological, well supported & \cite[ad loc. Mic 5:2]{keil-delitzsch-minor-prophets-1949}; \cite[ad loc. Matt 2:6]{france-matthew-2007} \\
\pvartarget{TestimonyArguments.BornInBethlehem}{7}$P_{7}$ & Jesus of Nazareth satisfies the Bethlehem criterion & interpretive & Christian, hist.-gramm., well supported & \cite[ad loc. Matt 2:6]{france-matthew-2007} \\
\bottomrule
\end{longtable}

\noindent\footnotesize \textsuperscript{\dag} grounded in scripture alone: the reading is assumed, not argued for.
\normalsize

\par\medskip\noindent\textbf{\texttt{quotation\allowbreak{}Edge}}\par\nobreak\noindent
Matthew explicitly quotes Micah — uncontroversial verbal citation, attested across traditions.

{\small\begin{verbatim}
def quotationEdge : IntertextEdge :=
  { fromPassage := Scripture.matthew2_6, toPassage := Scripture.micah5_2,
    relation := RelationType.quotation,
    source := Scripture.na28Apparatus Scripture.matthew2_6 }
\end{verbatim}}
\par\medskip\noindent\textbf{\texttt{predictive\allowbreak{}Reading}}\par\nobreak\noindent
The Christian predictive reading of Micah 5:2.

{\small\begin{verbatim}
def predictiveReading : Interpretation :=
  { passage := Scripture.micah5_2,
    reading := "Micah 5:2 predicts the Messiah's birthplace as Bethlehem",
    asRelation := some RelationType.prediction,
    source :=
      {
        primary :=
          Reference.work Bib.keilDelitzschMinorProphets
            (Bib.Locus.adLoc Scripture.micah5_2),
        tradition := Tradition.christianTypological,
        confidence := Confidence.wellSupported } }
\end{verbatim}}
\par\medskip\noindent\textbf{\texttt{born\allowbreak{}In\allowbreak{}Bethlehem}}\par\nobreak\noindent
The criterion this argument establishes a candidate must meet.

{\small\begin{verbatim}
def bornInBethlehem : FulfillmentCriterion :=
  { name := "born in Bethlehem", basis := predictiveReading }
\end{verbatim}}
\par\medskip\noindent\textbf{\texttt{cite}}\par\nobreak\noindent
Citation and classification for every atom. Total, so nothing is uncited.

{\small\begin{verbatim}
def cite : Claim → AtomMeta
\end{verbatim}}
\paragraph*{The line of reason}
\leavevmode\par\nobreak
One strand, and the argument’s whole weight is on it.

\par\medskip\noindent\textbf{\texttt{to\allowbreak{}Criterion}}\par\nobreak\noindent
From the predictive reading and Matthew’s intent, the criterion follows.

\begin{derivation}
  \item[] \((\pvar{TestimonyArguments.BornInBethlehem}{1} \land \pvar{TestimonyArguments.BornInBethlehem}{3}) \rightarrow \pvar{TestimonyArguments.BornInBethlehem}{6}\)\gloss{\emph{would give:} The Messiah must be born in Bethlehem}
\end{derivation}
\par\medskip\noindent\textbf{\texttt{to\allowbreak{}Fulfilment}}\par\nobreak\noindent
From the criterion and the historical claim, the fulfilment follows.

\begin{derivation}
  \item[] \((\pvar{TestimonyArguments.BornInBethlehem}{6} \land \pvar{TestimonyArguments.BornInBethlehem}{4}) \rightarrow \pvar{TestimonyArguments.BornInBethlehem}{7}\)\gloss{\emph{would give:} Jesus of Nazareth satisfies the Bethlehem criterion}
\end{derivation}
\par\medskip\noindent\textbf{\texttt{predictive\allowbreak{}Line}} --- Predictive reading of Micah 5:2\par\nobreak\noindent
\textbf{The predictive line.} Micah’s oracle read as prophecy, Matthew’s quotation read as a fulfilment claim, and the criterion that follows.

\begin{derivation}
  \item[(1)] \(\pvar{TestimonyArguments.BornInBethlehem}{1}\)\gloss{Micah 5:2 is a forward-looking Messianic prediction}
  \item[(2)] \(\pvar{TestimonyArguments.BornInBethlehem}{2}\)\gloss{Matthew 2:6 quotes Micah 5:2}
  \item[(3)] \(\pvar{TestimonyArguments.BornInBethlehem}{3}\)\gloss{Matthew’s quotation intends predictive fulfilment}
  \item[(4)] \((\pvar{TestimonyArguments.BornInBethlehem}{1} \land \pvar{TestimonyArguments.BornInBethlehem}{3}) \rightarrow \pvar{TestimonyArguments.BornInBethlehem}{6}\)
  \item[\(\vdash\)] \(\pvar{TestimonyArguments.BornInBethlehem}{6}\)\gloss{\emph{at issue:} The Messiah must be born in Bethlehem}
\end{derivation}
\par\medskip\noindent\textbf{\texttt{critical\allowbreak{}Line}} --- Critical reading of Micah 5:2 as a near-term oracle\par\nobreak\noindent
The critical line: the quotation is granted and the prediction denied, so the same step delivers nothing.

\begin{derivation}
  \item[(1)] \(\pvar{TestimonyArguments.BornInBethlehem}{2}\)\gloss{Matthew 2:6 quotes Micah 5:2}
  \item[(2)] \(\pvar{TestimonyArguments.BornInBethlehem}{5}\)\gloss{Micah 5:2 is a near-term oracle about a contemporary Judaean ruler}
  \item[(3)] \(\lnot \pvar{TestimonyArguments.BornInBethlehem}{1}\)\gloss{\emph{denied:} Micah 5:2 is a forward-looking Messianic prediction}
  \item[(4)] \((\pvar{TestimonyArguments.BornInBethlehem}{1} \land \pvar{TestimonyArguments.BornInBethlehem}{3}) \rightarrow \pvar{TestimonyArguments.BornInBethlehem}{6}\)
  \item[\(\vdash\)] \(\pvar{TestimonyArguments.BornInBethlehem}{6}\)\gloss{\emph{at issue:} The Messiah must be born in Bethlehem}
\end{derivation}
\paragraph*{Packages}
\leavevmode\par\nobreak
\par\medskip\noindent\textbf{\texttt{christian}} --- Christian predictive reading of Micah 5:2\par\nobreak\noindent
The Christian predictive argument.

\begin{derivation}
  \item[(1)] \(\pvar{TestimonyArguments.BornInBethlehem}{1}\)\gloss{Micah 5:2 is a forward-looking Messianic prediction}
  \item[(2)] \(\pvar{TestimonyArguments.BornInBethlehem}{2}\)\gloss{Matthew 2:6 quotes Micah 5:2}
  \item[(3)] \(\pvar{TestimonyArguments.BornInBethlehem}{3}\)\gloss{Matthew’s quotation intends predictive fulfilment}
  \breakrow{(4)} \(\pvar{TestimonyArguments.BornInBethlehem}{4}\)\gloss{Jesus of Nazareth was born in Bethlehem}
  \breakrow{(5)} \((\pvar{TestimonyArguments.BornInBethlehem}{1} \land \pvar{TestimonyArguments.BornInBethlehem}{3}) \rightarrow \pvar{TestimonyArguments.BornInBethlehem}{6}\)\gloss{\emph{would give:} The Messiah must be born in Bethlehem}
  \item[(6)] \((\pvar{TestimonyArguments.BornInBethlehem}{6} \land \pvar{TestimonyArguments.BornInBethlehem}{4}) \rightarrow \pvar{TestimonyArguments.BornInBethlehem}{7}\)
  \item[\(\vdash\)] \(\pvar{TestimonyArguments.BornInBethlehem}{7}\)\gloss{\emph{at issue:} Jesus of Nazareth satisfies the Bethlehem criterion}
  \keepwithnext
\end{derivation}
\emph{Grounded in scripture alone:} Jesus of Nazareth was born in Bethlehem.
\par\medskip\noindent\textbf{\texttt{critical}} --- Critical reading of Micah 5:2 as a near-term oracle\par\nobreak\noindent
The critical reading: the quotation is granted, the prediction is not.

\begin{derivation}
  \item[(1)] \(\pvar{TestimonyArguments.BornInBethlehem}{2}\)\gloss{Matthew 2:6 quotes Micah 5:2}
  \item[(2)] \(\pvar{TestimonyArguments.BornInBethlehem}{5}\)\gloss{Micah 5:2 is a near-term oracle about a contemporary Judaean ruler}
  \item[(3)] \(\lnot \pvar{TestimonyArguments.BornInBethlehem}{1}\)\gloss{\emph{denied:} Micah 5:2 is a forward-looking Messianic prediction}
  \breakrow{(4)} \((\pvar{TestimonyArguments.BornInBethlehem}{1} \land \pvar{TestimonyArguments.BornInBethlehem}{3}) \rightarrow \pvar{TestimonyArguments.BornInBethlehem}{6}\)\gloss{\emph{would give:} The Messiah must be born in Bethlehem}
  \item[(5)] \((\pvar{TestimonyArguments.BornInBethlehem}{6} \land \pvar{TestimonyArguments.BornInBethlehem}{4}) \rightarrow \pvar{TestimonyArguments.BornInBethlehem}{7}\)
  \item[\(\vdash\)] \(\pvar{TestimonyArguments.BornInBethlehem}{7}\)\gloss{\emph{at issue:} Jesus of Nazareth satisfies the Bethlehem criterion}
  \keepwithnext
\end{derivation}
\emph{Grounded in scripture alone:} Jesus of Nazareth was born in Bethlehem.
\paragraph*{Results}
\leavevmode\par\nobreak
\par\medskip\noindent\textbf{\texttt{christian\_\allowbreak{}establishes}}\par\nobreak\noindent
Given the Christian premises, the conclusion follows.

{\small\begin{verbatim}
theorem christian_establishes : Establishes christian
-- axioms: propext, Quot.sound
\end{verbatim}}
\par\medskip\noindent\textbf{\texttt{critical\allowbreak{}Reading}}\par\nobreak\noindent
The critical reading, written down: Micah’s oracle concerns a contemporary Judaean ruler, so no messianic birthplace criterion arises and nothing about Jesus follows from it.

{\small\begin{verbatim}
def criticalReading : Valuation Claim :=
  fun a =>
    match a with
    | Claim.micahPredictsBethlehem => False
    | Claim.messiahBornInBethlehem => False
    | Claim.jesusSatisfiesCriterion => False
    | x => True
\end{verbatim}}
\par\medskip\noindent\textbf{\texttt{critical\_\allowbreak{}not\_\allowbreak{}establishes}}\par\nobreak\noindent
The critical reading does not establish the conclusion.

{\small\begin{verbatim}
theorem critical_not_establishes : ¬Establishes critical
-- axioms: propext, Quot.sound
\end{verbatim}}
\par\medskip\noindent\textbf{\texttt{jesus\_\allowbreak{}satisfies\_\allowbreak{}bethlehem}}\par\nobreak\noindent
End to end: under the Christian package, Jesus satisfies the Bethlehem criterion. This is the first result in the library that runs from cited premises through a machine-checked entailment to a claim about a person.

{\small\begin{verbatim}
theorem jesus_satisfies_bethlehem : Satisfies jesus bornInBethlehem
-- axioms: propext, Quot.sound
\end{verbatim}}
\paragraph*{Satisfiability}
\leavevmode\par\nobreak
\texttt{Entails} is vacuously true over a premise set with no model, so a package built from contradictory premises would establish its conclusion with every gate passing. Only the package carrying a positive \texttt{Establishes} result needs checking: \texttt{critical} has a countermodel already, and a countermodel is a valuation satisfying every premise.

Not tagged \texttt{@[headline]} — a claim about the health of the encoding rather than about Micah.

\par\medskip\noindent\textbf{\texttt{everything\allowbreak{}Holds\allowbreak{}Reading}}\par\nobreak\noindent
The reading on which every claim in this argument holds at once.

{\small\begin{verbatim}
def everythingHoldsReading : Valuation Claim :=
  fun x => True
\end{verbatim}}
\par\medskip\noindent\textbf{\texttt{christian\_\allowbreak{}is\_\allowbreak{}satisfiable}}\par\nobreak\noindent
The Christian package has a model, so \texttt{christian\_\allowbreak{}establishes} is not vacuous.

{\small\begin{verbatim}
theorem christian_is_satisfiable : Satisfiable christian.premises
-- axioms: propext, Quot.sound
\end{verbatim}}

\bibliographystyle{plain}
\bibliography{references}
\end{document}
